\section{家庭、劳作与知识的成化正德证据}

本组材料把视线从朝廷决策移向家庭、劳作、环境与知识如何在文书中留下痕迹。土地、户籍、税役、河工、学校、宗教空间与地方灾情，往往只在征解、工程、诉讼、捐输或褒奖的片段中出现。材料的缺口并非无关紧要：它说明哪些人的行动更容易进入档案，哪些人的经验只能从地方记录和物质遗存中间接追索。

\subsection{山区家庭与编户}

\eventanchor{ML-1470-003}{明·1470（成化六年）}
\paragraph{1470（成化六年）：景泰、天顺时期的继承、人事与诏令链条可在本年材料中定位，政变责任不可由单一史书裁断} 这一锚点使家庭、劳动或地方资源进入可追查的历史时间。账本注明的把握范围是来源导向的年度桥接条目：本年材料可定位相关政务线索；具体月日、发文机关、地域和执行结果仍须逐项回查，不能以制度文本或奏报替代实际效果。《宪宗实录》成化六年条；《明史·本纪》及相关列传；诏令、奏疏或会典版本 提供了定位事件的材料链，却分别反映编纂者、报送者和保存者的视角。其发生条件可概括为继承、官僚协调或皇权运作的压力使问题进入朝廷程序；随后可见的制度或社会影响是它影响后续人事与政策议程；动机和派系标签不能由单一叙事裁定。对于日常生活、性别与家庭分工、非精英劳作或未留名者，仍不能由这些材料直接补足。桥接条目只标示可回查的年度问题与材料链；实录有中央选择性，崇祯期尤其需要奏疏、地方、朝鲜及满文材料交叉。

\eventanchor{ML-1470-004}{明·1470（成化六年）}
\paragraph{1470（成化六年）：蒙古诸部、朝鲜及西北绿洲的交涉在本年留下多方材料，应以具体使团和地名校正概称} 这一锚点使家庭、劳动或地方资源进入可追查的历史时间。账本注明的把握范围是来源导向的年度桥接条目：本年材料可定位相关政务线索；具体月日、发文机关、地域和执行结果仍须逐项回查，不能以制度文本或奏报替代实际效果。《宪宗实录》成化六年条；《明史·外国传》；《朝鲜王朝实录》、琉球《历代宝案》或海外档案 提供了定位事件的材料链，却分别反映编纂者、报送者和保存者的视角。其发生条件可概括为朝贡名义、边境安全与港口贸易的利益使交涉持续发生；随后可见的制度或社会影响是它为后续册封、互市或海防安排留下文书与跨政权比较线索。对于日常生活、性别与家庭分工、非精英劳作或未留名者，仍不能由这些材料直接补足。桥接条目只标示可回查的年度问题与材料链；实录有中央选择性，崇祯期尤其需要奏疏、地方、朝鲜及满文材料交叉。

\eventanchor{ML-1471-001}{明·1471（成化七年）}
\paragraph{1471（成化七年）：刘通的叛军被击败} 这一锚点使家庭、劳动或地方资源进入可追查的历史时间。账本注明的把握范围是编年候选的年序可由官修与后出材料交叉；具体月日、参与者、战果或数字必须回查所列材料，不能将单一叙事当作定论。《宪宗实录》成化七年条；《明史·兵志》及相关列传；边镇、地方志或对方材料 提供了定位事件的材料链，却分别反映编纂者、报送者和保存者的视角。其发生条件可概括为前期边防、地方武装或跨区域资源调动的压力推动了该行动；随后可见的制度或社会影响是它改变了后续调兵、补给或地方秩序的条件；战果和伤亡须分开核对。对于日常生活、性别与家庭分工、非精英劳作或未留名者，仍不能由这些材料直接补足。译写候选以编年材料定位；实录记载反映朝廷选择，崇祯期须另以奏疏、地方及敌对政权材料交叉。

\eventanchor{ML-1471-002}{明·1471（成化七年）}
\paragraph{1471（成化七年）：军饷、漕运和户籍赋役的本年记录揭示恢复秩序的行政成本，账面额与实收须分列。（材料核查重点：诏令或奏议的形成与发受链）} 这一锚点使家庭、劳动或地方资源进入可追查的历史时间。账本注明的把握范围是来源导向的年度桥接条目：本年材料可定位相关政务线索；具体月日、发文机关、地域和执行结果仍须逐项回查，不能以制度文本或奏报替代实际效果。《宪宗实录》成化七年条；《明会典》相关条目；地方志、黄册/契约/河工材料 提供了定位事件的材料链，却分别反映编纂者、报送者和保存者的视角。其发生条件可概括为财政征解、交通工程或货币支付的压力使相关措施进入政务；随后可见的制度或社会影响是其法定安排不等于地方执行，后续须以地域材料追踪负担与成效。对于日常生活、性别与家庭分工、非精英劳作或未留名者，仍不能由这些材料直接补足。桥接条目只标示可回查的年度问题与材料链；实录有中央选择性，崇祯期尤其需要奏疏、地方、朝鲜及满文材料交叉。；本条特别核对诏令或奏议的形成与发受链。

\eventanchor{ML-1471-003}{明·1471（成化七年）}
\paragraph{1471（成化七年）：军饷、漕运和户籍赋役的本年记录揭示恢复秩序的行政成本，账面额与实收须分列。（材料核查重点：报称信息的来源、抄传与行政分类）} 这一锚点使家庭、劳动或地方资源进入可追查的历史时间。账本注明的把握范围是来源导向的年度桥接条目：本年材料可定位相关政务线索；具体月日、发文机关、地域和执行结果仍须逐项回查，不能以制度文本或奏报替代实际效果。《宪宗实录》成化七年条；《明会典》相关条目；地方志、黄册/契约/河工材料 提供了定位事件的材料链，却分别反映编纂者、报送者和保存者的视角。其发生条件可概括为财政征解、交通工程或货币支付的压力使相关措施进入政务；随后可见的制度或社会影响是其法定安排不等于地方执行，后续须以地域材料追踪负担与成效。对于日常生活、性别与家庭分工、非精英劳作或未留名者，仍不能由这些材料直接补足。桥接条目只标示可回查的年度问题与材料链；实录有中央选择性，崇祯期尤其需要奏疏、地方、朝鲜及满文材料交叉。；本条特别核对报称信息的来源、抄传与行政分类。

\eventanchor{ML-1471-004}{明·1471（成化七年）}
\paragraph{1471（成化七年）：科举、学校和法律条文在本年制度材料中可追踪，不能据规范文本推定州县实践一致} 这一锚点使家庭、劳动或地方资源进入可追查的历史时间。账本注明的把握范围是来源导向的年度桥接条目：本年材料可定位相关政务线索；具体月日、发文机关、地域和执行结果仍须逐项回查，不能以制度文本或奏报替代实际效果。《宪宗实录》成化七年条；《明史·选举志》或《艺文志》；文集、碑刻或书目材料 提供了定位事件的材料链，却分别反映编纂者、报送者和保存者的视角。其发生条件可概括为制度、教育或知识传播的需要推动了相关行动或文本形成；随后可见的制度或社会影响是它提供制度和传播史的锚点，不能据此推断社会整体接受。对于日常生活、性别与家庭分工、非精英劳作或未留名者，仍不能由这些材料直接补足。桥接条目只标示可回查的年度问题与材料链；实录有中央选择性，崇祯期尤其需要奏疏、地方、朝鲜及满文材料交叉。

\subsection{劳役、市场与交通}

\eventanchor{ML-1472-001}{明·1472（成化八年）}
\paragraph{1472（成化八年）：蒙古诸部、朝鲜及西北绿洲的交涉在本年留下多方材料，应以具体使团和地名校正概称。（材料核查重点：诏令或奏议的形成与发受链）} 这一锚点使家庭、劳动或地方资源进入可追查的历史时间。账本注明的把握范围是来源导向的年度桥接条目：本年材料可定位相关政务线索；具体月日、发文机关、地域和执行结果仍须逐项回查，不能以制度文本或奏报替代实际效果。《宪宗实录》成化八年条；《明史·外国传》；《朝鲜王朝实录》、琉球《历代宝案》或海外档案 提供了定位事件的材料链，却分别反映编纂者、报送者和保存者的视角。其发生条件可概括为朝贡名义、边境安全与港口贸易的利益使交涉持续发生；随后可见的制度或社会影响是它为后续册封、互市或海防安排留下文书与跨政权比较线索。对于日常生活、性别与家庭分工、非精英劳作或未留名者，仍不能由这些材料直接补足。桥接条目只标示可回查的年度问题与材料链；实录有中央选择性，崇祯期尤其需要奏疏、地方、朝鲜及满文材料交叉。；本条特别核对诏令或奏议的形成与发受链。

\eventanchor{ML-1472-002}{明·1472（成化八年）}
\paragraph{1472（成化八年）：蒙古诸部、朝鲜及西北绿洲的交涉在本年留下多方材料，应以具体使团和地名校正概称。（材料核查重点：报称信息的来源、抄传与行政分类）} 这一锚点使家庭、劳动或地方资源进入可追查的历史时间。账本注明的把握范围是来源导向的年度桥接条目：本年材料可定位相关政务线索；具体月日、发文机关、地域和执行结果仍须逐项回查，不能以制度文本或奏报替代实际效果。《宪宗实录》成化八年条；《明史·外国传》；《朝鲜王朝实录》、琉球《历代宝案》或海外档案 提供了定位事件的材料链，却分别反映编纂者、报送者和保存者的视角。其发生条件可概括为朝贡名义、边境安全与港口贸易的利益使交涉持续发生；随后可见的制度或社会影响是它为后续册封、互市或海防安排留下文书与跨政权比较线索。对于日常生活、性别与家庭分工、非精英劳作或未留名者，仍不能由这些材料直接补足。桥接条目只标示可回查的年度问题与材料链；实录有中央选择性，崇祯期尤其需要奏疏、地方、朝鲜及满文材料交叉。；本条特别核对报称信息的来源、抄传与行政分类。

\eventanchor{ML-1472-003}{明·1472（成化八年）}
\paragraph{1472（成化八年）：地方灾异、赈济与水利条目为本年社会史提供入口，地方志的修纂层次需要说明} 这一锚点使家庭、劳动或地方资源进入可追查的历史时间。账本注明的把握范围是来源导向的年度桥接条目：本年材料可定位相关政务线索；具体月日、发文机关、地域和执行结果仍须逐项回查，不能以制度文本或奏报替代实际效果。《宪宗实录》成化八年条；《明史·五行志》；受灾府县地方志与奏疏 提供了定位事件的材料链，却分别反映编纂者、报送者和保存者的视角。其发生条件可概括为自然风险与地方报告促成朝廷或地方的应对记录；随后可见的制度或社会影响是赈济、迁徙和损失的实际范围仍须以受灾地区材料分别核验。对于日常生活、性别与家庭分工、非精英劳作或未留名者，仍不能由这些材料直接补足。桥接条目只标示可回查的年度问题与材料链；实录有中央选择性，崇祯期尤其需要奏疏、地方、朝鲜及满文材料交叉。

\eventanchor{MM-1472-TURPAN-HAMI-SEIZURE}{明·1472（成化八年）}
\paragraph{1472（成化八年）：吐鲁番阿黑麻势力进攻哈密，明朝在东天山的册封与朝贡网络受到冲击} 这一锚点使家庭、劳动或地方资源进入可追查的历史时间。账本注明的把握范围是进攻及明廷接获边报可靠；哈密城控制、当地王族选择和吐鲁番军队规模主要来自外部奏报，不能视为完整现场记录。《宪宗实录》成化八年哈密条；《明史·西域传》；《明实录》所载哈密使臣文书 提供了定位事件的材料链，却分别反映编纂者、报送者和保存者的视角。其发生条件可概括为哈密王族继承脆弱，而吐鲁番寻求控制绿洲、贸易和使团通道；明廷远距支援受嘉峪关外补给限制；随后可见的制度或社会影响是明廷须在册封、关市限制和边防成本之间重新取舍；1473年的继嗣交涉未消除吐鲁番压力。对于日常生活、性别与家庭分工、非精英劳作或未留名者，仍不能由这些材料直接补足。因此，登记、报送与实际经验之间仍须保留距离。

\paragraph{1473（成化九年）：明军与蒙古盟友联合进攻哈密，但在蒙古人放弃他们后撤退} 这一锚点使家庭、劳动或地方资源进入可追查的历史时间。账本注明的把握范围是编年候选的年序可由官修与后出材料交叉；具体月日、参与者、战果或数字必须回查所列材料，不能将单一叙事当作定论。《宪宗实录》成化九年条；《明史·兵志》及相关列传；边镇、地方志或对方材料 提供了定位事件的材料链，却分别反映编纂者、报送者和保存者的视角。其发生条件可概括为前期边防、地方武装或跨区域资源调动的压力推动了该行动；随后可见的制度或社会影响是它改变了后续调兵、补给或地方秩序的条件；战果和伤亡须分开核对。对于日常生活、性别与家庭分工、非精英劳作或未留名者，仍不能由这些材料直接补足。译写候选以编年材料定位；实录记载反映朝廷选择，崇祯期须另以奏疏、地方及敌对政权材料交叉。

\paragraph{1473（成化九年）：科举、学校和法律条文在本年制度材料中可追踪，不能据规范文本推定州县实践一致} 这一锚点使家庭、劳动或地方资源进入可追查的历史时间。账本注明的把握范围是来源导向的年度桥接条目：本年材料可定位相关政务线索；具体月日、发文机关、地域和执行结果仍须逐项回查，不能以制度文本或奏报替代实际效果。《宪宗实录》成化九年条；《明史·选举志》或《艺文志》；文集、碑刻或书目材料 提供了定位事件的材料链，却分别反映编纂者、报送者和保存者的视角。其发生条件可概括为制度、教育或知识传播的需要推动了相关行动或文本形成；随后可见的制度或社会影响是它提供制度和传播史的锚点，不能据此推断社会整体接受。对于日常生活、性别与家庭分工、非精英劳作或未留名者，仍不能由这些材料直接补足。桥接条目只标示可回查的年度问题与材料链；实录有中央选择性，崇祯期尤其需要奏疏、地方、朝鲜及满文材料交叉。

\subsection{灾异、工程与物质环境}

\paragraph{1473（成化九年）：土木危机后的边镇整顿、京师防务或复辟前后军政调度在本年材料中构成连续问题} 这一锚点使家庭、劳动或地方资源进入可追查的历史时间。账本注明的把握范围是来源导向的年度桥接条目：本年材料可定位相关政务线索；具体月日、发文机关、地域和执行结果仍须逐项回查，不能以制度文本或奏报替代实际效果。《宪宗实录》成化九年条；《明史·兵志》及相关列传；边镇、地方志或对方材料 提供了定位事件的材料链，却分别反映编纂者、报送者和保存者的视角。其发生条件可概括为前期边防、地方武装或跨区域资源调动的压力推动了该行动；随后可见的制度或社会影响是它改变了后续调兵、补给或地方秩序的条件；战果和伤亡须分开核对。对于日常生活、性别与家庭分工、非精英劳作或未留名者，仍不能由这些材料直接补足。桥接条目只标示可回查的年度问题与材料链；实录有中央选择性，崇祯期尤其需要奏疏、地方、朝鲜及满文材料交叉。

\paragraph{1473（成化九年）：哈密王位继承与朝贡使团问题引发明廷交涉，吐鲁番势力介入东天山政治} 这一锚点使家庭、劳动或地方资源进入可追查的历史时间。账本注明的把握范围是哈密政局和明廷交涉可见于实录；当地王位传承、部众立场和吐鲁番影响范围在中原材料中较难确证。《宪宗实录》成化九年哈密条；《明史·西域传》；《明实录》所载哈密使臣文书 提供了定位事件的材料链，却分别反映编纂者、报送者和保存者的视角。其发生条件可概括为哈密既是明朝西域朝贡体系的节点，也受本地王族和吐鲁番力量竞争牵动；随后可见的制度或社会影响是明廷不得不在册封、赏赐、关市与边防之间权衡；哈密问题此后反复成为西北边政负担。对于日常生活、性别与家庭分工、非精英劳作或未留名者，仍不能由这些材料直接补足。因此，登记、报送与实际经验之间仍须保留距离。

\paragraph{1474（成化十年）：于子军指挥重建和扩建长城，南封鄂尔多斯} 这一锚点使家庭、劳动或地方资源进入可追查的历史时间。账本注明的把握范围是编年候选的年序可由官修与后出材料交叉；具体月日、参与者、战果或数字必须回查所列材料，不能将单一叙事当作定论。《宪宗实录》成化十年条；《明史·本纪》及相关列传；诏令、奏疏或会典版本 提供了定位事件的材料链，却分别反映编纂者、报送者和保存者的视角。其发生条件可概括为继承、官僚协调或皇权运作的压力使问题进入朝廷程序；随后可见的制度或社会影响是它影响后续人事与政策议程；动机和派系标签不能由单一叙事裁定。对于日常生活、性别与家庭分工、非精英劳作或未留名者，仍不能由这些材料直接补足。译写候选以编年材料定位；实录记载反映朝廷选择，崇祯期须另以奏疏、地方及敌对政权材料交叉。

\paragraph{1474（成化十年）：科举、学校和法律条文在本年制度材料中可追踪，不能据规范文本推定州县实践一致} 这一锚点使家庭、劳动或地方资源进入可追查的历史时间。账本注明的把握范围是来源导向的年度桥接条目：本年材料可定位相关政务线索；具体月日、发文机关、地域和执行结果仍须逐项回查，不能以制度文本或奏报替代实际效果。《宪宗实录》成化十年条；《明史·选举志》或《艺文志》；文集、碑刻或书目材料 提供了定位事件的材料链，却分别反映编纂者、报送者和保存者的视角。其发生条件可概括为制度、教育或知识传播的需要推动了相关行动或文本形成；随后可见的制度或社会影响是它提供制度和传播史的锚点，不能据此推断社会整体接受。对于日常生活、性别与家庭分工、非精英劳作或未留名者，仍不能由这些材料直接补足。桥接条目只标示可回查的年度问题与材料链；实录有中央选择性，崇祯期尤其需要奏疏、地方、朝鲜及满文材料交叉。

\paragraph{1474（成化十年）：地方灾异、赈济与水利条目为本年社会史提供入口，地方志的修纂层次需要说明} 这一锚点使家庭、劳动或地方资源进入可追查的历史时间。账本注明的把握范围是来源导向的年度桥接条目：本年材料可定位相关政务线索；具体月日、发文机关、地域和执行结果仍须逐项回查，不能以制度文本或奏报替代实际效果。《宪宗实录》成化十年条；《明史·五行志》；受灾府县地方志与奏疏 提供了定位事件的材料链，却分别反映编纂者、报送者和保存者的视角。其发生条件可概括为自然风险与地方报告促成朝廷或地方的应对记录；随后可见的制度或社会影响是赈济、迁徙和损失的实际范围仍须以受灾地区材料分别核验。对于日常生活、性别与家庭分工、非精英劳作或未留名者，仍不能由这些材料直接补足。桥接条目只标示可回查的年度问题与材料链；实录有中央选择性，崇祯期尤其需要奏疏、地方、朝鲜及满文材料交叉。

\paragraph{1474（成化十年）：景泰、天顺时期的继承、人事与诏令链条可在本年材料中定位，政变责任不可由单一史书裁断} 这一锚点使家庭、劳动或地方资源进入可追查的历史时间。账本注明的把握范围是来源导向的年度桥接条目：本年材料可定位相关政务线索；具体月日、发文机关、地域和执行结果仍须逐项回查，不能以制度文本或奏报替代实际效果。《宪宗实录》成化十年条；《明史·本纪》及相关列传；诏令、奏疏或会典版本 提供了定位事件的材料链，却分别反映编纂者、报送者和保存者的视角。其发生条件可概括为继承、官僚协调或皇权运作的压力使问题进入朝廷程序；随后可见的制度或社会影响是它影响后续人事与政策议程；动机和派系标签不能由单一叙事裁定。对于日常生活、性别与家庭分工、非精英劳作或未留名者，仍不能由这些材料直接补足。桥接条目只标示可回查的年度问题与材料链；实录有中央选择性，崇祯期尤其需要奏疏、地方、朝鲜及满文材料交叉。

\subsection{教育、文书与地方知识}

\paragraph{1475（成化十一年）：湖南苗族叛乱遭镇压} 这一锚点使家庭、劳动或地方资源进入可追查的历史时间。账本注明的把握范围是编年候选的年序可由官修与后出材料交叉；具体月日、参与者、战果或数字必须回查所列材料，不能将单一叙事当作定论。《宪宗实录》成化十一年条；《明史·本纪》及相关列传；诏令、奏疏或会典版本 提供了定位事件的材料链，却分别反映编纂者、报送者和保存者的视角。其发生条件可概括为继承、官僚协调或皇权运作的压力使问题进入朝廷程序；随后可见的制度或社会影响是它影响后续人事与政策议程；动机和派系标签不能由单一叙事裁定。对于日常生活、性别与家庭分工、非精英劳作或未留名者，仍不能由这些材料直接补足。译写候选以编年材料定位；实录记载反映朝廷选择，崇祯期须另以奏疏、地方及敌对政权材料交叉。

\paragraph{1475（成化十一年）：学校、考试、刻书或礼制材料在本年留下文化制度线索，精英文本不能代表全体社会} 这一锚点使家庭、劳动或地方资源进入可追查的历史时间。账本注明的把握范围是来源导向的年度桥接条目：本年材料可定位相关政务线索；具体月日、发文机关、地域和执行结果仍须逐项回查，不能以制度文本或奏报替代实际效果。《宪宗实录》成化十一年条；《明史·选举志》或《艺文志》；文集、碑刻或书目材料 提供了定位事件的材料链，却分别反映编纂者、报送者和保存者的视角。其发生条件可概括为制度、教育或知识传播的需要推动了相关行动或文本形成；随后可见的制度或社会影响是它提供制度和传播史的锚点，不能据此推断社会整体接受。对于日常生活、性别与家庭分工、非精英劳作或未留名者，仍不能由这些材料直接补足。桥接条目只标示可回查的年度问题与材料链；实录有中央选择性，崇祯期尤其需要奏疏、地方、朝鲜及满文材料交叉。

\paragraph{1475（成化十一年）：本年灾荒、河工或赈恤的报奏材料可与地方志互校，报灾范围和实际损失应分别记录。（材料核查重点：诏令或奏议的形成与发受链）} 这一锚点使家庭、劳动或地方资源进入可追查的历史时间。账本注明的把握范围是来源导向的年度桥接条目：本年材料可定位相关政务线索；具体月日、发文机关、地域和执行结果仍须逐项回查，不能以制度文本或奏报替代实际效果。《宪宗实录》成化十一年条；《明史·五行志》；受灾府县地方志与奏疏 提供了定位事件的材料链，却分别反映编纂者、报送者和保存者的视角。其发生条件可概括为自然风险与地方报告促成朝廷或地方的应对记录；随后可见的制度或社会影响是赈济、迁徙和损失的实际范围仍须以受灾地区材料分别核验。对于日常生活、性别与家庭分工、非精英劳作或未留名者，仍不能由这些材料直接补足。桥接条目只标示可回查的年度问题与材料链；实录有中央选择性，崇祯期尤其需要奏疏、地方、朝鲜及满文材料交叉。；本条特别核对诏令或奏议的形成与发受链。

\paragraph{1475（成化十一年）：本年灾荒、河工或赈恤的报奏材料可与地方志互校，报灾范围和实际损失应分别记录。（材料核查重点：报称信息的来源、抄传与行政分类）} 这一锚点使家庭、劳动或地方资源进入可追查的历史时间。账本注明的把握范围是来源导向的年度桥接条目：本年材料可定位相关政务线索；具体月日、发文机关、地域和执行结果仍须逐项回查，不能以制度文本或奏报替代实际效果。《宪宗实录》成化十一年条；《明史·五行志》；受灾府县地方志与奏疏 提供了定位事件的材料链，却分别反映编纂者、报送者和保存者的视角。其发生条件可概括为自然风险与地方报告促成朝廷或地方的应对记录；随后可见的制度或社会影响是赈济、迁徙和损失的实际范围仍须以受灾地区材料分别核验。对于日常生活、性别与家庭分工、非精英劳作或未留名者，仍不能由这些材料直接补足。桥接条目只标示可回查的年度问题与材料链；实录有中央选择性，崇祯期尤其需要奏疏、地方、朝鲜及满文材料交叉。；本条特别核对报称信息的来源、抄传与行政分类。

\paragraph{1475（成化十一年）：本年灾荒、河工或赈恤的报奏材料可与地方志互校，报灾范围和实际损失应分别记录。（材料核查重点：同年地方材料所见的执行范围）} 这一锚点使家庭、劳动或地方资源进入可追查的历史时间。账本注明的把握范围是来源导向的年度桥接条目：本年材料可定位相关政务线索；具体月日、发文机关、地域和执行结果仍须逐项回查，不能以制度文本或奏报替代实际效果。《宪宗实录》成化十一年条；《明史·五行志》；受灾府县地方志与奏疏 提供了定位事件的材料链，却分别反映编纂者、报送者和保存者的视角。其发生条件可概括为自然风险与地方报告促成朝廷或地方的应对记录；随后可见的制度或社会影响是赈济、迁徙和损失的实际范围仍须以受灾地区材料分别核验。对于日常生活、性别与家庭分工、非精英劳作或未留名者，仍不能由这些材料直接补足。桥接条目只标示可回查的年度问题与材料链；实录有中央选择性，崇祯期尤其需要奏疏、地方、朝鲜及满文材料交叉。；本条特别核对同年地方材料所见的执行范围。

\paragraph{1476（成化十二年）六月：襄阳周边的流浪人口叛乱，直到政府允许他们减税索取土地} 这一锚点使家庭、劳动或地方资源进入可追查的历史时间。账本注明的把握范围是编年候选的年序可由官修与后出材料交叉；具体月日、参与者、战果或数字必须回查所列材料，不能将单一叙事当作定论。《宪宗实录》成化十二年条；《明会典》相关条目；地方志、黄册/契约/河工材料 提供了定位事件的材料链，却分别反映编纂者、报送者和保存者的视角。其发生条件可概括为财政征解、交通工程或货币支付的压力使相关措施进入政务；随后可见的制度或社会影响是其法定安排不等于地方执行，后续须以地域材料追踪负担与成效。对于日常生活、性别与家庭分工、非精英劳作或未留名者，仍不能由这些材料直接补足。译写候选以编年材料定位；实录记载反映朝廷选择，崇祯期须另以奏疏、地方及敌对政权材料交叉。
